\section{Estimation par filtre de Kalman}
\label{sec:kalman}

La méthode des moments n'utilise que les variations quotidiennes du spread. Elle ignore la dynamique de son
niveau, donc la mesure historique, et traite chaque série comme une observation exacte du facteur. Or le facteur
$y$ n'est pas observé : seuls le sont des rendements de différentes maturités, qui en dépendent avec une erreur.
Le cadre naturel est alors celui des modèles espace-état, estimés par maximum de vraisemblance à l'aide du filtre
de Kalman \citep[cours~5, section~I.2]{MPR2025}. Le filtre estime ensemble les paramètres risque-neutres, qui
fixent les chargements des rendements sur les facteurs, et les paramètres historiques, qui fixent la dynamique
des facteurs ; il reconstruit aussi la trajectoire des facteurs. On l'emploie ici en contrôle de la GMM, et pour
tester ce que les données disent de la mesure historique.

\subsection{Représentation espace-état}

On observe chaque jour ouvré $t$ les taux de swaps €STR $R_t(\tau)$ et les spreads OAT–€STR $S_t(\tau)$ aux six
tenors $\tau \in \{5, 10, 15, 20, 25, 30\}$ ans, soit un vecteur $z_t$ de $m = 12$ séries, centrées sur leur moyenne
empirique. L'état est le couple non observé $w_t = (x_t, y_t)'$. Dans le modèle, le taux zéro €STR de maturité
$\tau$ est affine en $x$, de pente $H_{a_x^Q}(\tau)/\tau$, et le spread zéro est affine en $y$, de pente
$H_{a_y^Q}(\tau)/\tau$. D'où l'équation de mesure (cours~5, définition~I.1)
\begin{equation}
z_t = B\,w_t + \varepsilon_t,\qquad
B = \begin{pmatrix} \big(H_{a_x^Q}(\tau_j)/\tau_j\big)_j & 0 \\ 0 & \big(H_{a_y^Q}(\tau_j)/\tau_j\big)_j \end{pmatrix},
\end{equation}
avec $\varepsilon_t \sim \mathcal N(0,\Omega)$ et $\Omega = \mathrm{diag}(\omega_1^2, \dots, \omega_{12}^2)$,
où les chargements dépendent des vitesses de retour risque-neutres $a^Q$, celles qui figurent dans les prix des
zéro-coupons. Sous la mesure historique, les facteurs suivent des processus d'Ornstein-Uhlenbeck de vitesses
$a^P$, avec les mêmes volatilités et la même corrélation ; les dynamiques sous les deux mesures sont spécifiées
séparément, comme dans la modélisation mixte du cours~1 (slide~132). Discrétisés exactement au pas
$\Delta = 1/252$, ils donnent l'équation de transition, un VAR(1) gaussien :
\begin{equation}
\begin{split}
w_t &= \Phi\,w_{t-1} + \eta_t,\qquad \Phi = \mathrm{diag}\big(e^{-a_x^P\Delta},\,e^{-a_y^P\Delta}\big),\\
\Var(\eta_t) &= \begin{pmatrix} \sigma_x^2\,I(2a_x^P) & \rho\sigma_x\sigma_y\,I(a_x^P + a_y^P) \\
\rho\sigma_x\sigma_y\,I(a_x^P + a_y^P) & \sigma_y^2\,I(2a_y^P)\end{pmatrix},
\end{split}
\end{equation}
avec $I(a) = (1-e^{-a\Delta})/a$. Les paramètres sont
$\theta = (a_x^Q, a_y^Q, a_x^P, a_y^P, \sigma_x, \sigma_y, \rho, \omega_1, \dots, \omega_{12})$.

C'est la représentation espace-état de Diebold, Rudebusch et Aruoba \citep[cours~3, slides~119-120]{MPR2025},
avec des chargements contraints par l'absence d'arbitrage : $H_a(\tau)/\tau$ est le chargement de pente de
Nelson-Siegel, qui tend vers le chargement de niveau, égal à 1, quand $a \to 0$. Une vitesse risque-neutre nulle
fait donc du facteur un facteur de niveau.

\subsection{Vraisemblance et identification}

Le modèle est linéaire et gaussien : le filtre de Kalman donne exactement la loi de $w_t$ sachant les observations
passées (cours~5, proposition~I.2), et la vraisemblance se calcule à partir des erreurs de prévision
$\nu_t = z_t - B\,w_{t|t-1}$ et de leurs variances $S_t$ :
\begin{equation}
\log L(\theta) = -\frac{mT}{2}\log 2\pi - \frac12\sum_{t=1}^T\Big(\log|S_t| + \nu_t'\,S_t^{-1}\,\nu_t\Big).
\end{equation}
Le filtre est initialisé à la loi stationnaire de l'état, et les quelques valeurs manquantes des séries €STR sont
traitées en réduisant ces jours-là la dimension de $z_t$ (cours~5, remarques~I.2 et~I.3). La vraisemblance est
maximisée par l'algorithme L-BFGS ; le filtre, la vraisemblance et le lissage sont ceux de la classe
\texttt{MLEModel} de statsmodels.

Laissée libre, la vraisemblance est maximale lorsque l'une des variances d'erreur s'annule : le facteur passe alors
exactement par une série, et la vraisemblance a plusieurs maxima selon la série choisie. On fixe ce choix en
supposant les deux taux à vingt ans, tenor de l'OAT sous-jacente, observés sans erreur. C'est l'hypothèse des
variables parfaitement modélisées de la technique d'inversion de \citet{ChenScott1993}
(cours~5, section~I.3, hypothèse~I.1) : les facteurs se lisent sur les taux à vingt ans, et les dix autres séries,
mesurées avec erreur, identifient les vitesses risque-neutres par la forme de leurs chargements.

\subsection{Ce que le filtre peut identifier : test sur données simulées}

Avant l'estimation sur données réelles, on vérifie ce que l'estimateur retrouve quand le modèle est vrai. On simule
48 échantillons de 1\,460 jours, la taille des données, aux paramètres estimés plus bas, à trois exceptions près :
$a_y^Q = 0{,}05$, pour savoir si une vitesse risque-neutre modeste serait détectée, et $a_x^P = a_y^P = 0{,}2$,
ordre de grandeur courant d'une vitesse historique (demi-vie de 3,5 ans). Le modèle est réestimé sur chaque
échantillon.

\begin{table}[htbp]
\centering\small
\caption{Estimation sur 48 échantillons simulés (volatilités en \%)}
\label{tab:kalman_sim}
\begin{tabular}{@{}lccccc@{}}
\toprule
 & Valeur vraie & Médiane & Écart-type & Quantile 5~\% & Quantile 95~\% \\
\midrule
$a_x^Q$ & 0,0118 & 0,0117 & 0,0006 & 0,0109 & 0,0126 \\
$a_y^Q$ & 0,0500 & 0,0499 & 0,0015 & 0,0480 & 0,0526 \\
$a_x^P$ & 0,2000 & 0,8173 & 0,6231 & 0,2291 & 2,3023 \\
$a_y^P$ & 0,2000 & 0,6166 & 0,7056 & 0,1042 & 1,7754 \\
$\sigma_x$ & 0,7751 & 0,7734 & 0,0163 & 0,7439 & 0,7970 \\
$\sigma_y$ & 0,3313 & 0,3301 & 0,0081 & 0,3183 & 0,3420 \\
$\rho$ & 0,1513 & 0,1511 & 0,0253 & 0,1147 & 0,1930 \\
\bottomrule
\end{tabular}
\end{table}

Les paramètres qui fixent les chargements et les covariances sont retrouvés avec précision
(tableau~\ref{tab:kalman_sim}) : une vitesse risque-neutre de 0,05 serait détectée sans ambiguïté. Les vitesses
historiques ne le sont pas. Sur six ans de données, un processus de demi-vie 3,5 ans ne parcourt qu'environ un
cycle, et la vraisemblance distingue mal une vitesse de 0,2 d'une vitesse de 1. L'estimateur est de plus biaisé :
en petit échantillon, le coefficient autorégressif $\phi = e^{-a^P\Delta}$ d'un processus persistant est
sous-estimé d'environ $(1+3\phi)/T$ (approximation de Kendall, cours~5, slides~52-53), donc la vitesse
$a^P \approx (1-\phi)/\Delta$ est surestimée d'environ $4/(T\Delta)$, soit 0,7 pour $T\Delta \approx 5{,}8$ années.
C'est l'ordre de grandeur de l'écart entre les médianes estimées, 0,82 et 0,62, et la vraie valeur. On retrouve le
problème de persistance du cours (section~I.6).

\subsection{Résultats sur données françaises}

\begin{table}[htbp]
\centering\small
\caption{Estimation par filtre de Kalman, 1\,460 jours du 4 janvier 2021 au 8 septembre 2026 (volatilités en \%)}
\label{tab:kalman}
\begin{tabular}{@{}lccc@{}}
\toprule
 & Filtre de Kalman & Écart-type & GMM ou swaptions \\
\midrule
$a_x^Q$ & 0,0118 & 0,0004 & 0,0107 \\
$a_y^Q$ & 0,0001 & 0,0004 & 0,001 (borne) \\
$a_x^P$ & 0,1216 & 0,1587 & — \\
$a_y^P$ & 0,2152 & 0,3591 & — \\
$\sigma_x$ & 0,7751 & 0,0123 & 0,7367 \\
$\sigma_y$ & 0,3313 & 0,0038 & 0,3374 \\
$\rho$ & 0,1513 & 0,0211 & 0,1299 \\
\bottomrule
\end{tabular}
\end{table}

Les estimations confirment la mesure de pricing (tableau~\ref{tab:kalman}). La vitesse risque-neutre du spread est
nulle, comme dans la GMM ; la volatilité du spread, 33,1~pb par an, et sa corrélation avec les taux, 0,15, sont
proches de celles de la GMM. La vitesse risque-neutre des taux, tirée de la seule coupe de la courbe €STR, vaut
0,0118, proche de celle qu'impliquent les swaptions (0,0107). Les vitesses historiques ne sont pas significatives,
ce que le test sur données simulées laissait prévoir.

Les écarts-types des erreurs de mesure vont de 5 à 27~pb pour les swaps €STR et de 4 à 12~pb pour les spreads,
les plus élevés au tenor de cinq ans. Le test de \citet{LjungBox1978}, à dix retards, rejette l'indépendance des
erreurs de prévision pour chaque série (p-valeur au plus $2{,}9\times10^{-261}$) : les écarts de chaque tenor au
facteur sont persistants. Un facteur par courbe ne reproduit pas les déformations de pente, que la deuxième
composante principale du spread mettait déjà en évidence ; la figure~\ref{fig:kalman} le montre au tenor de cinq
ans. C'est la conclusion du test $J$ de la GMM.

\begin{figure}[htbp]
\centering
\includegraphics[width=\textwidth]{kalman}
\caption{À gauche, vitesses historiques estimées sur les 48 échantillons simulés (valeur vraie 0,2) ; à droite,
spreads observés et spreads du modèle à 5 et 30 ans, avec le facteur lissé.}
\label{fig:kalman}
\end{figure}

Le filtre de Kalman confirme ainsi l'hypothèse~H2 pour la mesure de pricing et la rejette pour la mesure
historique. La volatilité, la corrélation et la vitesse risque-neutre du spread sont identifiées et concordent
avec la GMM ; six ans de données ne suffisent pas, en revanche, à identifier sa dynamique historique. La GMM reste
l'estimateur retenu, et la section~\ref{sec:resultats} mesure l'effet sur le prix d'un passage aux paramètres du
filtre.
